\documentclass[11pt]{article}

\usepackage[margin=1in]{geometry}
\usepackage{amsmath,amssymb,amsthm}
\usepackage{graphicx}
\usepackage[hidelinks]{hyperref}

\newtheorem{theorem}{Theorem}

\newcommand{\N}{\mathbb{N}}
\newcommand{\HFC}{\mathrm{HFC}}
\newcommand{\Lip}{\operatorname{Lip}}
\newcommand{\diam}{\operatorname{diam}}

\title{Literature-Implied Partial Subcase for arXiv:2004.04805}
\author{FA/Banach run fa\_banach\_001}
\date{June 14, 2026}

\begin{document}
\maketitle

\section*{Status}

This is a \texttt{literature\_implied\_answer (partial subcase)} packet. It
does not claim a new solution to the equi-coarse embedding problem in
Baudier--Lancien--Motakis--Schlumprecht \cite{BLMS2020}. It records that a
later theorem of Fovelle \cite{Fovelle2021} rules out the stronger
equi-Lipschitz embeddings of the Hamming graphs into the particular space
$T^*(T^*)$.

\section*{Source Problem}

Baudier--Lancien--Motakis--Schlumprecht ask whether the Hamming graphs fail to
equi-coarsely embed into every reflexive asymptotic-subsequential-$c_0$ space.
They single out the space $T^*(T^*)$:
\[
  \text{do the Hamming graphs fail to equi-coarsely embed into } T^*(T^*)?
\]
They had already proved that, if such embeddings exist, they cannot be of the
canonical array/tree forms described earlier in their paper.

\begin{figure}[h]
\centering
\includegraphics[width=.95\textwidth]{figures/open_problem_crop.png}
\caption{Source evidence from arXiv:2004.04805, page 26: final remarks and
Problem 5.1 asking about equi-coarse embeddings of Hamming graphs into
reflexive asymptotic-subsequential-$c_0$ spaces and into $T^*(T^*)$.}
\end{figure}

\section*{Supporting Literature}

Fovelle later proves that $T^*(T^*)$ does not have property
$\mathrm{HC}_\infty$, but does have property $\HFC_{p,d}$ for every
$p\in(1,\infty)$. The same paper notes that these finite-$p$ concentration
properties prevent equi-Lipschitz embeddings of the Hamming graphs.

\begin{figure}[h]
\centering
\includegraphics[width=.95\textwidth]{figures/supporting_theorem_crop.png}
\caption{Supporting evidence from arXiv:2106.04297, page 17: the corollary
stating that $T^*(T^*)$ has property $\HFC_{p,d}$ for every
$p\in(1,\infty)$.}
\end{figure}

\section*{Definitions and Identification}

Write $H_k^\omega=([\N]^k,d_H)$ for the $k$-dimensional Hamming graph. A
family $(H_k^\omega)_k$ equi-Lipschitz embeds into a Banach space $Y$ if
there are constants $A,B>0$ and maps $f_k:H_k^\omega\to Y$ such that
\[
  A\,d_H(\bar m,\bar n)\le \|f_k(\bar m)-f_k(\bar n)\|
  \le B\,d_H(\bar m,\bar n)
\]
for all $k$ and all $\bar m,\bar n\in[\N]^k$.

For $p\in(1,\infty)$, property $\HFC_{p,d}$ says that there is a constant
$\lambda$ such that, for every $k$ and every Lipschitz map
$f:([\N]^k,d_H)\to Y$, there is an infinite $\mathbb M\subseteq\N$ satisfying
\[
  \|f(\bar m)-f(\bar n)\|
  \le \lambda \left(\sum_{j=1}^k \alpha_j^p\right)^{1/p}
\]
for all $\bar m,\bar n\in[\mathbb M]^k$, where $\alpha_j$ is the directional
Lipschitz constant obtained by changing only the $j$th coordinate.

\section*{Idea of the Proof}

The source problem is about coarse embeddings, whose compression may grow
sublinearly. Fovelle's finite-$p$ concentration theorem is strong enough to
rule out the linear-scale case. Indeed, an equi-Lipschitz embedding would
force disjoint $k$-sets to be separated by order $k$, while $\HFC_{p,d}$
forces a large sub-Hamming graph to have diameter at most order $k^{1/p}$.
Since $p>1$, these estimates are incompatible for large $k$.

\begin{theorem}
The Hamming graphs $(H_k^\omega)_{k\ge 1}$ do not equi-Lipschitz embed into
$T^*(T^*)$.
\end{theorem}

\begin{proof}
Choose any $p\in(1,\infty)$. By Fovelle's corollary
\cite{Fovelle2021}, $T^*(T^*)$ has property $\HFC_{p,d}$; let $\lambda$ be a
constant witnessing this property.

Assume, for contradiction, that there are constants $A,B>0$ and maps
$f_k:H_k^\omega\to T^*(T^*)$ such that
\[
  A\,d_H(\bar m,\bar n)\le \|f_k(\bar m)-f_k(\bar n)\|
  \le B\,d_H(\bar m,\bar n)
\]
for every $k$ and all $\bar m,\bar n\in[\N]^k$.

For each $k$, every directional Lipschitz constant $\alpha_j$ of $f_k$ is at
most $B$, because changing one coordinate changes the Hamming distance by
one. Applying $\HFC_{p,d}$ to $f_k$, we obtain an infinite
$\mathbb M\subseteq\N$ such that
\[
  \|f_k(\bar m)-f_k(\bar n)\|
  \le \lambda\left(\sum_{j=1}^k B^p\right)^{1/p}
  = \lambda B k^{1/p}
\]
for all $\bar m,\bar n\in[\mathbb M]^k$.

Since $\mathbb M$ is infinite, choose disjoint $\bar m,\bar n\in[\mathbb M]^k$.
Then $d_H(\bar m,\bar n)=k$, so the lower Lipschitz estimate gives
\[
  A k \le \|f_k(\bar m)-f_k(\bar n)\| \le \lambda B k^{1/p}.
\]
This is impossible for all sufficiently large $k$, because $p>1$. Hence no
such equi-Lipschitz family exists.
\end{proof}

\section*{Verification and Limitations}

The packet uses no computational verifier. The proof is a direct consequence
of Fovelle's concentration property and the definitions of equi-Lipschitz
embedding and Hamming distance.

The limitation is essential: this proves only the failure of
equi-Lipschitz embeddings into $T^*(T^*)$. It does not rule out equi-coarse
embeddings with sublinear compression, which is the question asked in
\cite{BLMS2020}.

\section*{Novelty and Search Bounds}

This is not an originality claim. Before packaging, the run's cheap indexes
were searched for \texttt{2004.04805}, Hamming graphs, Johnson graphs,
$T^*(T^*)$, asymptotic-subsequential-$c_0$, and $\HFC$ keywords. No prior
packet or attempt for this exact source problem or subcase was found. Local
corpus and web searches for the Hamming/Johnson and $T^*(T^*)$ phrases found
the source paper and later concentration-property papers, with
arXiv:2106.04297 providing the supporting theorem used here.

\section*{Human Review Recommendation}

Review as a provenance/status packet. The main points to check are that
Fovelle's $\HFC_{p,d}$ convention uses the same Hamming graph metric as the
source paper and that the finite-$p$ concentration argument above correctly
rules out equi-Lipschitz embeddings. Do not count this as a full answer to the
source paper's equi-coarse problem.

\begin{thebibliography}{9}

\bibitem{BLMS2020}
F. Baudier, G. Lancien, P. Motakis, and Th. Schlumprecht,
\emph{The geometry of Hamming-type metrics and their embeddings into Banach
spaces},
arXiv:2004.04805.

\bibitem{Fovelle2021}
A. Fovelle,
\emph{Hamming graphs and concentration properties in non-quasi-reflexive
Banach spaces},
arXiv:2106.04297.

\end{thebibliography}

\end{document}
